\textbf{COLUNI 2026 -- Questão 29}\par\medskip Um estudo internacional publicado em 2025 por cientistas da Noruega, Austrália e Suécia, com quase 40 mil adolescentes, mostrou que cada hora extra de uso de telas antes de dormir pode aumentar em 59\% o risco de insônia e reduzir, em média, 24 minutos do tempo de sono por noite. Esses efeitos estão diretamente relacionados à redução da atividade da glândula pineal, causada pela exposição à luz azul das telas, que prejudica a produção da melatonina, o hormônio regulador do sono.

\begin{figure}[htbp]
\centering
\includegraphics[width=0.52\linewidth]{static/simulados/coluni/figuras/2026/questao-29/figura-1.png}
\end{figure}

Disponível em: https://www.cnnbrasil.com.br/saude/fica-no-celular-antes-de-dormir-entenda-como-o-habito-prejudica-seu-sono/. Acesso em: 21 ago. 2025.

Considerando a glândula pineal, é CORRETO afirmar que se trata de uma glândula:\begin{enumerate}
\item endócrina, que possui ducto e libera hormônio em cavidades do corpo.
\item endócrina, que não possui ducto e libera hormônio na corrente sanguínea.
\item exócrina, que possui ducto e libera hormônio na superfície do corpo.
\item exócrina, que não possui ducto e libera hormônios nas próprias células que os produziram.\end{enumerate}